\section{Motivation}
% Budget: 1 page.

Emerging equity markets are exposed to two kinds of shocks. Some are global: changes in US monetary policy, in the dollar or in investors' appetite for risk move all of them at once, a pattern that \citet{rey2013dilemma} describes as a global financial cycle. Others are local, such as an election, a currency crisis or a change in economic policy in one country, and the question is whether, and through which countries, they spread to the rest. Telling the two apart matters for several reasons. For investors, the benefit of diversifying across emerging markets depends on how much of their co-movement is a common component that no portfolio of these markets can diversify away. For policy makers, a high degree of connectedness calls for different responses if it reflects bilateral transmission (contagion) or a common exposure to the same external shock. For the empirical literature, the distinction is an old one: \citet{forbes2002contagion} showed that much apparent contagion during crises is high interdependence driven by common shocks and higher volatility.

The tools to measure connectedness have multiplied, and they do not all measure the same thing. The variance-decomposition approach of \citet{diebold2014network} summarizes how much of each market's forecast-error variance is due to shocks elsewhere; it is dense by construction and does not separate common from bilateral transmission. Sparse network estimators such as NETS \citep{barigozzi2019nets} keep only the links supported by the data, and FNETS \citep{barigozzi2024fnets} first removes the common factor and then estimates the network of what is left. Applied to the same data, these methods can give quite different pictures of how connected markets are, who transmits and who receives, and how stable the network is. Before using a connectedness measure to draw conclusions about contagion, it is therefore worth knowing how much of the measured connectedness depends on the estimation method.

Our sample of 17 emerging markets offers clear examples of both kinds of shock. In the week of the Argentine primary elections of August 2019 the Argentine fund (ARGT) lost 28.4\% (in log terms), its largest weekly fall in the sample: a large, country-specific shock. Seven months later, the COVID-19 crash of March 2020 produced the largest weekly losses of almost every market at the same time: 14 of the 17 markets recorded their worst week in March 2020. Argentina is one case among the 17 markets in this paper, but its 2019 episode illustrates the question we ask: whether a large local shock spreads to other markets, and whether the methods agree on that.
